% concatenated from TightTreesArxiv.tex
\documentclass[11pt]{article}

\usepackage[margin=1in]{geometry}
\usepackage[T1]{fontenc}
\usepackage{lmodern,comment}
\usepackage{amsmath,amssymb,amsthm,mathtools}
\usepackage{microtype}
\usepackage{enumitem}
\usepackage[hidelinks]{hyperref}
\usepackage[numbers,sort&compress]{natbib}
\newtheorem{theorem}{Theorem}[section]
\newtheorem{lemma}[theorem]{Lemma}
\newtheorem{corollary}[theorem]{Corollary}
\newtheorem{conjecture}{Conjecture}
\theoremstyle{remark}
\newtheorem{remark}[theorem]{Remark}
\numberwithin{equation}{section}

\DeclareMathOperator{\ex}{ex}
\DeclareMathOperator{\Perm}{Perm}
\DeclareMathOperator{\supp}{supp}
\newcommand{\sh}{\partial}
\newcommand{\Wcal}{\mathcal W}
\newcommand{\Scal}{\mathcal S}
\setlength{\parindent}{0pt}
\parskip=8pt
\hypersetup{
	pdftitle={Kalai's Conjecture for tight trees},
	pdfauthor={Dhruv Mubayi},
	pdfsubject={A permutation proof of the shadow bound for tight trees}
}

\title{Kalai's Conjecture for Tight Trees}
\author{
	Dhruv Mubayi\thanks{Department of Mathematics, Statistics and Computer Science, University of Illinois, Chicago, IL 60607. Email: mubayi@uic.edu. Research partially supported by NSF Awards DMS-2552740 and DMS-2153576.} \and
	Jacques Verstra\" ete\thanks{Department of Mathematics, University of California, San Diego, CA, 92093-0112 USA.
		Email: jverstraete@ucsd.edu.
		Research supported by NSF-BSF award DMS-2347832.  }}
\date{}

\begin{document}
	\maketitle
	
	\begin{abstract}
		Let $r \ge 2$ and $t \ge 1$. It is shown that if $T$ is an $r$-uniform tight tree with $t$ edges and
		$H$ is a $T$-free $r$-uniform hypergraph, then
		$|E(H)|\le (t-1)|\sh H|/r$, where $\sh H$ is the $(r-1)$-shadow
		of $H$. \iffalse Equality holds only for $(n,t + r - 2,r)$-designs.\fi The bound is tight infinitely often, and establishes Kalai's Conjecture, whose $r=2$ case is the Erd\H os-S\'os Conjecture.
		
		The proof was
		found by GPT-6 Astra, extending its method of proof for the Erd\H os-S\'os conjecture to the hypergraph setting. It is noteworthy that previous proofs of special cases of the Erd\H os-S\'os conjecture do not extend to give tights bounds in the hypergraph setting. 
		
		A strengthening of the Erd\H os-S\'os conjecture due to the authors about tight lower bounds on the number  of copies of a tree in a graph with  average degree $d \ge t-1\ge 0$ remains open. 
	\end{abstract}
	
	\section{AI Declaration Statement}
	The proof of the main result was found by GPT-6 Astra by first  prompting it to extend its recent proof of the Erd\H os-S\'os conjecture (see Adamczewski and Bloom~\cite[Appendix B.4]{Epoch})  to the case of hypergraph tight paths as in~\cite{FJKMVpaths}, and then prompting it to prove the full Kalai conjecture. As we elaborate throughout the text, GPT-6 Astra's  proof technique for Erd\H os-S\'os conjecture, which uses permutations of the underlying vertex set, perhaps has its genesis in an old argument of Perles (see~\cite{KupitzPerles}), which considered convex geometric graphs (these corresponds to cyclic permutations of the vertex set).  This method was developed in the hypergraph case by F\"uredi-Jiang-Kostochka-Mubayi-Verstra\"ete in~\cite{FJKMVpaths} and gave the previously best bounds for Kalai's conjecture for tight paths. The authors have checked the proof and rewritten much of the text provided by GPT-6 Astra and added further explanations throughout the paper. They take full responsibility for all the content in the paper.
	
	
	\section{Introduction}\label{sec:introduction}
	
	All hypergraphs in this paper are finite and simple. An $r$-uniform
	hypergraph, or \emph{$r$-graph}, has edges of size $r$. For an $r$-graph
	$T$, let $\ex(n,T)$ denote the maximum number of edges in an $n$-vertex
	$r$-graph containing no copy of $T$. 
	
	The Erd\H{o}s--Gallai Theorem~\cite{EG} asserts that a graph with no
	path of $t$ edges has at most $(t-1)n/2$ edges. The
	Erd\H{o}s--S\'os Conjecture~\cite{Erdos} asks for the same bound when
	the forbidden path is replaced by an arbitrary tree with $t$ edges. This conjecture has received extensive attention in the literature~\cite{ASS,BPSS2019,BPSS2021,BD,DPRS,Erdos,Erdos1964,EG,ErdosSos1970,FanSun2007,GoerlichZak2016,HRSW2020,McLennan2005,Pokrovskiy2024,Rozhon2019,ReedStein2026,SacleWozniak1997,Stein,TinerTomlin2022,Wozniak1996}. The Erd\H{o}s--S\'os Conjecture was recently proved employing a permutation argument by GPT-6 Astra AI, described by Adamczewski and Bloom~\cite[Appendix B.4]{Epoch}. Similar counting arguments are key to the method of Perles~\cite{KupitzPerles} giving a very elegant alternative proof of the Erd\H{o}s--Gallai Theorem using convex geometric graphs, and to extensions to uniform hypergraphs in~\cite{FJKMVpaths}.
	
	
	Kalai proposed a hypergraph generalization of the Erd\H{o}s--S\'{o}s Conjecture in 1984, recorded by
	Frankl and F\"uredi~\cite{FF}, using the following notion of a tree. An $r$-graph $T$ is a \emph{tight $r$-tree} if its edges have an ordering
	$e_1,\ldots,e_t$ such that for every $i>1$, there exists
	$z_i\in e_i$ and an index $p(i)<i$ satisfying
	\begin{equation}\label{eq:tight-tree}
		z_i\notin\bigcup_{h<i}e_h,
		\qquad e_i\setminus\{z_i\}\subseteq e_{p(i)}.
	\end{equation}
	We take $V(T)=\bigcup_{e\in E(T)}e$, so a tight $r$-tree with $t$
	edges has exactly $t+r-1$ vertices -- the case $r = 2$ is precisely a tree in the graph sense. Kalai's Conjecture is as follows:
	
	\begin{conjecture}\label{conj:kalai} {\bf (Kalai's Conjecture)} 
		For $n \geq r \geq 2$ and every tight $r$-tree $T$ with $t$ edges,
		\begin{equation}\label{eq:kalai-binomial}
			\ex(n,T)\le\frac{t-1}{r}\binom{n}{r-1}.
		\end{equation}
	\end{conjecture} 
	
	To motivate this conjecture, let $b = t + r - 2$ and take a family of $b$-vertex blocks in which each $(r-1)$-set
	lies in at most one block and almost every $(r-1)$-set is covered,
	and put a complete $r$-graph on each block. Any copy of a tight tree
	must stay in a single block -- two edges from different blocks cannot
	meet in $r-1$ vertices -- and therefore each tight tree has at most $b = t + r - 2$ vertices and therefore fewer than $t$ edges. Using
	R\"odl's packing theorem~\cite{Rodl} providing the existence of asymptotically optimal partial Steiner systems, one obtains as $n \rightarrow \infty$:
	\[
	\ex(n,T)\ge\left(\frac{t-1}{r}-o(1)\right) \cdot \binom{n}{r-1}.
	\]
	Keevash's existence theorem for designs~\cite{Keevash} (see also Glock, K\"{u}hn, Lo and Osthus~\cite{GKLO}) gives exact
	equality in~\eqref{eq:kalai-binomial} for infinitely many $n$ by
	covering every $(r-1)$-set exactly once with $b$-sets -- these are Steiner systems $S(r-1,b,n)$. In the special case that the tree $T$ is a sunflower -- every edge of $T$ has the form $R \cup \{v\}$ where $|R| = r - 1$ -- then any $r$-uniform hypergraph in which every $(r - 1)$-set is contained in exactly $t - 1$ hyperedges is $T$-free. These are sometimes called \emph{simple $(r - 1)$-$(n,r,t-1)$-designs}.
	
	
	Several cases of Kalai's Conjecture were previously known. Frankl
	and F\"uredi~\cite{FF} proved (\ref{eq:kalai-binomial}) for star-shaped tight
	trees, which have an edge meeting every other edge in $r-1$ vertices.
	F\"uredi, Jiang, Kostochka and the authors~\cite{FJKMVtrunk}
	proved an asymptotic version for tight trees with a bounded trunk and
	the exact bound for tight trees with at most four edges. Here a trunk
	is an initial tight subtree in a construction~\eqref{eq:tight-tree}
	such that each remaining edge attaches directly to an edge of that
	subtree; the asymptotic statement fixes $r$ and the trunk size and
	lets $t$ tend to infinity. Stein~\cite{Stein} proved (\ref{eq:kalai-binomial})
	for all tight trees when the host $r$-graph is $r$-partite.
	
	\subsection{Main Theorem}
	
	The \emph{shadow} of an $r$-graph $H$ is
	\[
	\sh H=\{e\subseteq V(H):|e|=r-1,\ e\subseteq h
	\text{ for some }h\in E(H)\}.
	\]
	We prove  Kalai's
	conjecture by proving the following shadow bound  (the equivalence to Kalai's Conjecture was established by F\"uredi, Jiang, Kostochka and the authors~\cite[Proposition 2.2]{FJKMVtrunk}).
	
	\begin{theorem}\label{thm:main}
		Let $n \ge r\ge 2$, and let $T$ be a tight $r$-tree with $t\ge1$ edges\iffalse, and $b = t + r - 2$\fi.
		Every $n$-vertex $T$-free $r$-graph $H$ satisfies
		\begin{equation}\label{eq:main-shadow}
			|E(H)| \; \le \; \frac{t-1}{r} \cdot |\sh H| \; \le \; \frac{t-1}{r} \cdot \binom{n}{r-1} .
		\end{equation}
		
	\end{theorem}
	
	
	This theorem also settles Conjecture 1.1 of  F\"uredi, Jiang, Kostochka and the authors~\cite{FJKMVpaths}. The proof is given in Section \ref{sec:trees}. 
	
	GPT-6 Astra has also give a proof that equality holds in Theorem~\ref{thm:main} if and only if $H$ is a Steiner $S(r-1,b,n)$-system ($b=t+r-2$) or $T$ is a sunflower and every $(r - 1)$-subset of $V(H)$ is contained in exactly $t - 1$ edges of $H$. We have not verified this claim regarding the equality characterization. 
	
	\subsection{Counting conjecture}
	
	The following generalization of the Erd\H{o}s-S\'{o}s Conjecture was stated by the authors in~\cite{MubayiV2016}: let $T$ be a tree with $t$ edges. Then there exists $d_0=d_0(t)$ such that every $n$-vertex graph $G$ of average degree $d\geq d_0$ contains at least \begin{equation}\label{eq:treecount} N_T(G) \geq n(d)_t = n d(d-1)(d-2)\cdots(d-t+1) \end{equation} labeled copies of $T$, where $N_T(G)$ denotes the number of injective graph homomorphisms from $T$ to $G$. 
	The authors in~\cite{MubayiV2016} proved the asymptotically sharp estimate \begin{equation} 
		N_T(G) \geq \left(1-O\left(\frac{t^5}{d^2}\right)\right)n(d)_t. \end{equation} 
	The conjectured bound (\ref{eq:treecount}) has subsequently been proved by Wilson~\cite{Wilson2025} for $d_0=\Omega(t^4)$. A natural strengthening posed by the authors would be to prove the result with \begin{equation} d_0(t) = t - 1. \end{equation} Such a statement would clearly imply the Erd\H{o}s--S\'os conjecture.  It is not too difficult to verify that the conjecture holds for certain small trees, as well as for double stars.
	A proof, which we have not verified, is claimed by GPT-6 Astra when $d=\Omega(t^3)$ for all trees $T$ and when $d=\Omega(t^2)$ and $T$ is a  path with $t$ edges. 
	
	
	
	We make the following 
	even more general hypergraph counting conjecture:
	
	\begin{conjecture} \label{conj:hypcount}
		Let $n \geq r \geq 2$, let $d \geq t - 1 \geq 0$, and let $H$ be an $n$-vertex $r$-uniform hypergraph with at least $\frac{d}{r}{n \choose r - 1}$ edges. Then for any tight tree $T$ with $t$ edges, 
		\begin{equation}
			N_T(H) \geq (r-1)!{n\choose r-1}(d)_t.
		\end{equation}    
	\end{conjecture}
	\bigskip
	
	In other words, any hypergraph of average $(r - 1)$-degree $d \geq t - 1$ contains at least $(r-1)!{n\choose r-1}(d)_t$ labelled copies of any prescribed tight tree $T$ with $t$ edges. This implies Kalai's Conjecture (Conjecture \ref{conj:kalai}), but, as mentioned above, remains open even for paths in graphs. GPT-6 Astra has claimed a proof of Conjecture~\ref{conj:hypcount} for $d>16r^2t^3$. We have not verified this. 
	
	\subsection{Organization} 
	
	The special case of trees that are {\em tight paths} contains the main ideas of the proof of Theorem \ref{thm:main}, and is presented in Section \ref{sec:paths}. 
	This proof generalizes the ideas in~\cite{FJKMVpaths}, which previously had the best bound for the extremal number of tight paths. The proof of Theorem \ref{thm:main} is given in subsequent sections.
	\iffalse
	and, in Section \ref{sec:equality}, we characterize the case of equality.
	\fi
	
	\section{Proof of Theorem \ref{thm:main} for tight paths}\label{sec:paths}
	
	The \emph{tight path} $P_k^r$ has vertex set
	$\{x_1,\ldots,x_{k+r-1}\}$ and edges
	$\{x_i,\ldots,x_{i+r-1}\}$ for $1\le i\le k$. Using convex geometric hypergraphs, F\"uredi, Jiang, Kostochka and the authors~\cite[Theorem 1.2]{FJKMVpaths} proved
	\begin{equation}\label{eq:previous-path}
		\ex(n,P_k^r)\le
		\begin{cases}
			\displaystyle\frac{k-1}{2}\binom{n}{r-1},&r\text{ even},\\[6pt]
			\displaystyle\frac12\left(k+\left\lfloor\frac{k-1}{r}\right\rfloor\right)
			\binom{n}{r-1},&r\text{ odd}.
		\end{cases}
	\end{equation}
	
	The proof of this theorem, which proceeds by considering a count of ordered tight paths in so-called convex geometric hypergraphs~\cite{FJKMVpaths} partly motivates the proof of Theorem~\ref{thm:main}; indeed in this section we prove the tight path case of Theorem \ref{thm:main}, as it is much simpler and yet illustrates the main ideas. Roughly speaking, we count permutations whose first $r-1$ vertices, together with a marked later vertex, form an edge of the host hypergraph. For paths, a cyclic rotation extends every eligible state except the first one in each permutation, and these ideas were used to study the extremal function for tight paths in convex geometric hypergraphs~\cite{FJKMVpaths}, and we give this simpler ``path argument'' proof of Theorem \ref{thm:main} in this section. The more complicated case of general trees is achieved in Sections \ref{subsec:rooted-states} -- \ref{subsec:operationsroot}.
	
	
	\subsection{Definitions and notation} \label{subsection:defns} Fix an $r$-graph $H$ on a vertex set $V$ of size $n\ge r$.
	For a permutation $\pi=(v_1,\ldots,v_n)$ of $V$, write
	\[
	R(\pi)=(v_1,\ldots,v_{r-1}),\qquad
	V_j(\pi)=\{v_1,\ldots,v_j\}.
	\]
	Let $\Wcal$ be the set of permutations whose first $r-1$ vertices
	form a member of $\sh H$, and define the set of \emph{marked states}
	by
	\begin{equation}\label{eq:states}
		\Omega=\bigl\{(\pi,j):\pi\in\Wcal,\ r\le j\le n,
		\ \{v_1,\ldots,v_{r-1},v_j\}\in E(H)\bigr\}.
	\end{equation}
	Set $W=|\Wcal|$ and $M=|\Omega|$. Direct counting gives
	\begin{equation}\label{eq:counts}
		\begin{aligned}
			W&=(r-1)!\,|\sh H|\,(n-r+1)!,\\
			M&=r!\,|E(H)|\,(n-r+1)!.
		\end{aligned}
	\end{equation}
	For the first identity, choose the initial shadow edge, order it,
	and order the remaining vertices. For the second, choose an edge,
	choose its marked vertex in $r$ ways, and order its other $r-1$
	vertices at the beginning of the permutation. Order all remaining
	$n-r+1$ vertices arbitrarily; the position of the chosen marked
	vertex determines $j$. For $\ell\ge1$, let $\Scal_\ell\subseteq\Omega$ consist of the
	states $(\pi,j)$ for which $H[V_j(\pi)]$ contains a tight path of
	$\ell$ edges ending with the ordered tuple $R(\pi)$. Explicitly,
	there must be distinct vertices $x_1,\ldots,x_{\ell+r-1}\in V_j(\pi)$
	such that
	\begin{equation}\label{eq:path-state}
		\mbox{For }1 \leq i \leq \ell, 
		\{x_i,\ldots,x_{i+r-1}\} \in E(H)\mbox{ and }
		(x_{\ell+1},\ldots,x_{\ell+r-1}) = R(\pi).
	\end{equation}
	The permutation specifies the allowed vertex set and the terminal
	tuple; the path may traverse its other vertices in any order.
	Every marked state supports the one-edge path
	$(v_j,v_1,\ldots,v_{r-1})$, so
	\begin{equation}\label{eq:path-base}
		\Scal_1=\Omega.
	\end{equation}
	
	\subsection{Mapping lemma for tight paths}
	
	\begin{lemma}\label{lem:path-transfer}
		For every $\ell\ge1$,
		\begin{equation}\label{eq:path-transfer}
			|\Scal_\ell|\le |\Scal_{\ell+1}|+W.
		\end{equation}
	\end{lemma}
	
	\begin{proof}
		For each $\pi\in\Wcal$ having a state in $\Scal_\ell$, discard
		$(\pi,j)$ with $j$ least among its states in $\Scal_\ell$.
		At most $W$ states are discarded. We provide an injection $\Phi'$ from the remaining states
		into $\Scal_{\ell+1}$.
		
		For any $(\pi,j)\in\Omega$, write
		\[
		\pi=(a_1,\ldots,a_{r-1},A,x,B),\qquad x=v_j,
		\]
		where $A=(v_r,\ldots,v_{j-1})$ and $B=(v_{j+1},\ldots,v_n)$
		are possibly empty words. Define
		\begin{equation}\label{eq:path-rotation}
			\Phi(\pi,j)=(\pi',j),\qquad
			\pi'=(a_2,\ldots,a_{r-1},x,A,a_1,B).
		\end{equation}
		This rotates the entries in positions $1,\ldots,r-1,j$ and fixes
		all other entries. It preserves the marked edge, and hence maps
		$\Omega$ to itself. Since $\Phi^r$ is the identity, $\Phi$ is a
		bijection of $\Omega$. Moreover,
		\begin{equation}\label{eq:path-support}
			V_j(\pi')=V_j(\pi),\qquad
			R(\pi')=(a_2,\ldots,a_{r-1},x).
		\end{equation}
		
		Now let $(\pi,j)\in\Scal_\ell$ be a state that has not been discarded in the first step. There exists
		$i<j$ with $(\pi,i)\in\Scal_\ell$. A witnessing path for this
		earlier state lies in $V_i(\pi)$ and ends with
		$(a_1,\ldots,a_{r-1})$. Since $x=v_j\notin V_i(\pi)$, appending
		$x$ introduces a new vertex and the new edge
		$\{a_1,\ldots,a_{r-1},x\}\in E(H)$.
		The resulting tight path has $\ell+1$ edges, lies in $V_j(\pi')$,
		and ends with $R(\pi')$. Thus $\Phi(\pi,j)\in\Scal_{\ell+1}$.
		The restriction of the bijection $\Phi$ is the required injection $\Phi'$.
	\end{proof}
	
	\begin{proof}[Proof of Theorem~\ref{thm:main} for tight paths]
		We aim to show that if $H$ is an $n$-vertex $r$-uniform $P_k^r$-free hypergraph, then $|E(H)| \leq (k - 1)/r \cdot |\sh H|$. The assertion is immediate if $E(H)=\varnothing$. Otherwise
		$n\ge r$, so the preceding definitions apply. If $H$ is
		$P_k^r$-free, then $\Scal_k=\varnothing$. Iterating
		Lemma~\ref{lem:path-transfer} and using~\eqref{eq:path-base} gives
		\[
		M=|\Scal_1|\le (k-1)W.
		\]
		Substituting~\eqref{eq:counts} and cancelling the factor
		$(r-1)!(n-r+1)!$ yields $r|E(H)|\le(k-1)|\sh H|$. \end{proof}
	
	\section{Proof of Theorem \ref{thm:main}}\label{sec:trees}
	
	For general tight trees, we prescribe the image of an entire root edge as in Adamczewski and Bloom~\cite[Appendix B.4]{Epoch}. A prefix permutation inequality then allows two rooted pieces to be joined
	without an additional error term. This part is not necessary in the proof given in Section \ref{sec:paths}, as we may repeatedly remove leaf edges and map them to leaf edges of shorter tight paths, whereas for a tight tree, we cannot map leaf edges to leaf edges.
	Induction then combines this joining
	operation with the addition of a leaf edge.
	
	
	We retain the notation $H$, $V$, $\Wcal$, $\Omega$, $W$ and $M$ from
	Section~\ref{sec:paths}. The proof for trees uses embeddings with a
	prescribed image of a whole root edge. 
	The following straightforward decomposition lemma allows us to do induction on the size of the tree.
	
	\begin{lemma}\label{lem:structure}
		Let $T$ be a tight $r$-tree with at least two edges, and fix
		$e\in E(T)$. At least one of the following holds.
		\begin{enumerate}[label=\textup{(\roman*)},leftmargin=*]
			\item\label{case:leaf}
			$T$ is obtained from a smaller tight tree $S$ by adding a new
			vertex $z$ and the edge $e=F\cup\{z\}$, where $F$ is an
			$(r-1)$-subset of an edge $f$ of $S$.
			\item\label{case:gluing}
			$T=T_1\cup T_2$ for tight trees $T_1,T_2$, each with fewer edges
			than $T$, such that
			\[
			E(T_1)\cap E(T_2)=\{e\},\qquad V(T_1)\cap V(T_2)=e.
			\]
		\end{enumerate}
	\end{lemma}
	
	This lemma is necessary for our inductive proof, which maps leaves to leaves when $T$ is a tight path (as happens in Case (i)), but which may map leaves to non-leaf edges when $T$ is not a path (Case (ii)). We omit the proof of Lemma~\ref{lem:structure}, which is an easy exercise.
	
	\iffalse \begin{proof}
		Choose an ordering and parents as in~\eqref{eq:tight-tree}, and
		let $D$ be the graph on $E(T)$ with edges $e_i e_{p(i)}$ for
		$i>1$. Each new node is attached to one earlier node, so $D$ is
		a graph tree. Adjacent nodes represent hyperedges meeting in
		exactly $r-1$ vertices.
		
		For each $v\in V(T)$, the nodes of $D$ whose hyperedges contain
		$v$ form a connected subtree. Indeed, the first occurrence of
		$v$ starts with one node. Each later hyperedge containing $v$
		attaches to a parent also containing $v$, because $v$ is then
		one of its old vertices.
		
		Every connected subtree of $D$ represents a tight subtree of
		$T$. To check this, root that graph subtree at any node and list
		its nodes with parents preceding children. Let a child hyperedge
		$g$ have parent $h$, and let $z$ be the unique vertex of
		$g\setminus h$. No earlier node is a descendant of $g$, so the
		path from $g$ to any earlier node passes through $h$. If an
		earlier hyperedge contained $z$, connectedness of the nodes
		containing $z$ would force $z\in h$, a contradiction. Thus $z$
		is new at this step, and the ordering is a tight-tree construction.
		
		If $e$ is a leaf of $D$, let $f$ be its neighbor and let $z$ be
		the unique vertex of $e\setminus f$. Connectedness of the nodes
		containing $z$ implies that no other edge contains $z$. Deleting
		$e$ and $z$ leaves the tight tree represented by $D-e$, proving
		\ref{case:leaf} with $F=e\cap f$.
		
		Otherwise $e$ has degree at least two in $D$. Partition the
		components of $D-e$ into two nonempty groups and add the node
		$e$ to each group. The resulting connected subtrees represent
		tight trees $T_1,T_2$, both smaller than $T$, whose edge sets
		intersect only in $e$. If a vertex outside $e$ occurred in both
		trees, the path between two nodes containing it would pass
		through $e$, contradicting connectedness of its occurrence nodes.
		Thus $V(T_1)\cap V(T_2)=e$, proving~\ref{case:gluing}.
	\end{proof}
	\fi
	\subsection{States rooted at an edge}\label{subsec:rooted-states}
	We now generalize the definition of ${\mathcal S}_{\ell}$ in Section~\ref{subsection:defns} from tight paths to arbitrary hypergraphs.
	Let $S$ be an $r$-graph with a distinguished ordered edge
	$\vec e=(u_1,\ldots,u_r)$. For $V(H)=\{v_1, \ldots, v_n\}$ and $\pi=v_1v_2\ldots v_n$, a state $(\pi,j)\in\Omega$ is
	\emph{good for $(S,\vec e)$} if there is an injective embedding
	$\varphi:V(S)\longrightarrow V(H)$ such that for $1 \leq i \leq r-1$,
	\begin{equation}\label{eq:rooted-state}
		\begin{gathered}
			\varphi(u_i)=v_i,\qquad \varphi(u_r)=v_j,\\
			\varphi(V(S))\subseteq V_j(\pi).
		\end{gathered}
	\end{equation}
	Here an embedding sends every edge of $S$ to an edge of $H$.
	Let $C(S,\vec e)$ be the number of good states. Each state is
	counted once, regardless of how many embeddings witness it.
	
	This number does not depend on the ordering of $e$. To see this,
	write $q_i=i$ for $i<r$ and $q_r=j$. If the root is reordered as
	$(u_{\sigma(1)},\ldots,u_{\sigma(r)})$, permute the host entries
	in these positions so that
	\[
	v'_{q_i}=v_{q_{\sigma(i)}}\qquad(1\le i\le r),
	\]
	and keep all other entries fixed. This is a bijection on $\Omega$,
	preserves $V_j(\pi)$, and makes the same embedding a witness for
	the reordered root. We therefore write $C(S,e)$, choosing an
	ordering of the fixed edge $e$ when convenient. The count can
	depend on the choice of $e$. For the one-edge tree,
	\begin{equation}\label{eq:rooted-base}
		C(e,e)=M.
	\end{equation}
	
	\subsection{Adding a leaf}\label{subsec:operationsleaf}
	
	The following lemma deals with Case (i) of Lemma~\ref{lem:structure} in the decomposition of a tight tree, and is the analog of the case of paths (Lemma~\ref{lem:path-transfer}) in Section \ref{sec:paths}.
	
	\begin{lemma}\label{lem:leaf}
		Let $S$ be an $r$-graph, let $f\in E(S)$, and let $F\subset f$
		have size $r-1$. Form $T$ by adding a new vertex $z$ and the
		edge $e=F\cup\{z\}$. Then
		\begin{equation}\label{eq:leaf}
			C(S,f)\le C(T,e)+W.
		\end{equation}
	\end{lemma}
	
	\begin{proof}
		Write $F=\{u_1,\ldots,u_{r-1}\}$ and $f=F\cup\{y\}$.
		Choose the root orders $(u_1,\ldots,u_{r-1},y)$ for $S$ and
		$(u_1,\ldots,u_{r-1},z)$ for $T$.
		For every $\pi\in\Wcal$, discard its first good state for
		$(S,f)$, if one exists. This discards at most $W$ states.
		
		
		For a retained state $(\pi,j)$, there is an earlier good state
		$(\pi,i)$ with $i<j$. Choose its witnessing embedding $\varphi$
		of $S$. It maps $u_h$ to $v_h$ for $h<r$, maps $y$ to $v_i$,
		and has image in $V_i(\pi)$. Since $v_j\notin V_i(\pi)$, extend
		$\varphi$ by sending $z$ to $v_j$. The additional edge maps to
		$\{v_1,\ldots,v_{r-1},v_j\}\in E(H)$, and the resulting embedding
		satisfies~\eqref{eq:rooted-state} for $(T,e)$ at $(\pi,j)$.
		Thus every retained state is itself good for $(T,e)$, proving
		\eqref{eq:leaf}.
	\end{proof}
	
	
	\subsection{Joining along root edges}\label{subsec:operationsroot}
	
	
	
	%\subsection{A permutation-prefix inequality}\label{subsec:prefix}
	
	The more complicated decomposition offered by Case (ii) of Lemma~\ref{lem:structure} is the subject of this section.
	For a finite set $U$, let $\Perm(U)$ denote the set of
	permutations of $U$. For $w=(w_1,\ldots,w_m)\in\Perm(U)$,
	write
	\[
	w_{\le s}=(w_1,\ldots,w_s),
	\qquad
	w_{>s}=(w_{s+1},\ldots,w_m).
	\]
	The support of a permutation or one of its blocks is
	the set of its entries. When $|U|=m$ and $\mathcal F\subseteq2^U$, put
	\[
	C_U(\mathcal F)=\{(w,s):w\in\Perm(U),\ 0\le s\le m,
	\ \supp(w_{\le s})\in\mathcal F\}\]  and  $c_U({\cal F})=|C_U({\cal F})|$.
	The empty prefix is included in $C_U(\mathcal F)$. Thus the total number of pairs
	$(w,s)$ is $(m+1)m!=(m+1)!$. 
	
	In Lemma~\ref{lem:prefix} below, we should think of a tight tree $T$ being expressed as the union of two tight trees $T_1$ and $T_2$ with a common edge $e$ so  $T=T_1 \cup T_2$ and $V(T_1) \cap V(T_2)=e$, 
	and $\cal A$, $\cal B$, $\cal C$ as vertex sets of the underlying hypergraph that support $T_1, T_2$, and $T$, respectively, where (the image of) $e$ is excluded.
	\begin{lemma}\label{lem:prefix}
		Suppose $\mathcal A,\mathcal B,\mathcal C\subseteq2^U$ satisfy
		$\mathcal C\subseteq\mathcal A$ and
		\begin{equation}\label{eq:disjoint-gluing}
			R\in\mathcal A,\quad X\in\mathcal B,\quad R\cap X=\varnothing
			\quad\Longrightarrow\quad R\cup X\in\mathcal C.
		\end{equation}
		If $|U|=m$, then
		\begin{equation}\label{eq:prefix-gluing}
			c_U(\mathcal A)+c_U(\mathcal B)
			\le (m+1)!+c_U(\mathcal C).
		\end{equation}
	\end{lemma}
	
	\begin{proof}
		We construct an injection from the set of pairs $(w,s)\in  
		C_U(\mathcal A)$  for which $\supp(w_{\le s})\not \in\mathcal C$ to the set of all pairs $(w', s')$ with  $\supp(w'_{\le s'})\not \in\mathcal B$. 
		The will imply
		$c_U(\mathcal A)-c_U(\mathcal C)
		\le (m+1)!-c_U(\mathcal B)$ as required.
		
		Let $(w,s)$ be a pair of the first
		kind, and let $R$ be the shortest prefix of $w$ with
		$\supp(R)\in\mathcal A\setminus\mathcal C$.
		This prefix exists and has length at most $s$. Write
		\[
		w=RXY,\qquad s=|R|+|X|,\qquad w'=XRY, \qquad s'=|X|
		\]
		where juxtaposition denotes concatenation, and define
		\begin{equation}\label{eq:prefix-map}
			(w,\,s)\longmapsto(w',\,s').
		\end{equation}
		If $\supp(X)\in\mathcal B$, then~\eqref{eq:disjoint-gluing}
		would imply $\supp(RX)\in\mathcal C$, contrary to the assumption
		on the input pair. Hence the output has prefix outside
		$\mathcal B$.
		
		
		Given $(w', s')$ let us now show how to recover $(w,s)$ uniquely, thus showing that this map is injective. First observe that, given $w'$, the integer  $s'$ determines the word
		$X$. In the remaining suffix $w'_{>s'}=RY$, take the shortest prefix whose
		support belongs to $\mathcal A \setminus \mathcal C$. This prefix
		is exactly $R$, since its support belongs to  $A\setminus\mathcal C$, and each of its shorter prefixes is
		the corresponding shorter prefix of the original word $w$, which
		does not belong to $\mathcal A \setminus \mathcal C$. We thereby recover $R$ and $Y$, and hence $(w,s)$. This also applies when $R$ or $X$ is empty.
	\end{proof}
	In the application to tight trees, we have the stronger
	inclusion $\mathcal C\subseteq\mathcal A\cap\mathcal B$:
	an embedding of the whole tree restricts to an embedding
	of either rooted subtree. However, the lemma requires
	only $\mathcal C\subseteq\mathcal A$.
	
	
	
	
	
	\begin{lemma}\label{lem:root-gluing}
		Suppose $T=T_1\cup T_2$, where
		\[
		V(T_1)\cap V(T_2)=e,\qquad e\in E(T_1)\cap E(T_2).
		\]
		Then
		\begin{equation}\label{eq:root-gluing}
			C(T_1,e)+C(T_2,e)\le M+C(T,e).
		\end{equation}
	\end{lemma}
	
	\begin{proof}
		Fix an ordering $e=(u_1,\ldots,u_r)$ and an assignment
		$u_i\mapsto a_i$ to distinct vertices with
		$\{a_1,\ldots,a_r\}\in E(H)$. Put
		\[
		A_0=\{a_1,\ldots,a_r\},\qquad U=V(H)\setminus A_0,
		\qquad m=n-r.
		\]
		A marked state with this root-edge assignment is of the form 
		$$((a_1, \ldots, a_{r-1}, P, a_r, Q), r+|P|).$$
		Consequently, for $w=(P, Q)$, these marked states are in bijection
		with all pairs $(w,s)$, where $w\in\Perm(U)$ and $0\le s\le m$:
		\begin{equation}\label{eq:fiber}
			(w,s)\longleftrightarrow
			\bigl((a_1,\ldots,a_{r-1},w_{\le s},a_r,w_{>s}),\ r+s\bigr).
		\end{equation}
		There are $(m+1)!$ such states. In particular, $s=0$ is allowed:
		the empty prefix outside $A_0$ corresponds to a marked state.
		
		
		Let $\mathcal A$, $\mathcal B$ and $\mathcal C$ be the families of
		sets $X\subseteq U$ for which $T_1$, $T_2$ and $T$, respectively,
		embed into $H[A_0\cup X]$ with the prescribed assignment on $e$.
		Clearly $\mathcal C\subseteq\mathcal A$, since if we restrict an embedding of $\mathcal C$ we get an embedding of $\mathcal A$.
		If $R\in\mathcal A$, $X\in\mathcal B$ and $R\cap X=\varnothing$,
		choose the two witnessing embeddings. They agree on every vertex
		of $e$ and have disjoint images outside $e$. Since
		$V(T_1)\cap V(T_2)=e$, their union is an injective embedding of
		$T$ into $H[A_0\cup R\cup X]$. Hence $R\cup X\in\mathcal C$.
		Consequently, $\mathcal A, \mathcal B, \mathcal C$ satisfy the hypothesis of Lemma~\ref{lem:prefix}.
		
		We now sum the conclusion of Lemma~\ref{lem:prefix} to obtain the conclusion of Lemma~\ref{lem:root-gluing}. For clarity of presentation, we make the counting explicit. Keep the ordering
		$e=(u_1,\ldots,u_r)$ fixed, and let
		$\mathcal I_H$ be the set of ordered edges of $H$, so that
		$|\mathcal I_H|=r!|E(H)|$.
		For each $\mathbf a=(a_1,\ldots,a_r)\in\mathcal I_H$, define
		\[
		\Omega_{\mathbf a}
		=
		\{(\pi,j)\in\Omega:
		v_i=a_i\text{ for }1\le i<r,\quad v_j=a_r\}.
		\]
		Let $C_{\mathbf a}(S,e)$ be the number of states in
		$\Omega_{\mathbf a}$ that are good for $(S,e)$.
		Every state $(\pi,j)\in\Omega$ determines the unique assignment
		$\mathbf a=(v_1,\ldots,v_{r-1},v_j)$.
		Consequently, the sets $\Omega_{\mathbf a}$ partition $\Omega$,
		and, for every $S$ under consideration,
		\begin{equation}\label{eq:assignment-sum}
			C(S,e)
			=
			\sum_{\mathbf a\in\mathcal I_H}C_{\mathbf a}(S,e).
		\end{equation}
		We count each good state once, regardless of how many embeddings satisfy its defining conditions.
		Now fix $\mathbf a\in\mathcal I_H$, and define $A_0$, $U$ and
		the families $\mathcal A,\mathcal B,\mathcal C$ as above.
		Under the bijection~\eqref{eq:fiber}, a pair $(w,s)$ corresponds
		to the state
		\[
		\pi=(a_1,\ldots,a_{r-1},
		w_1,\ldots,w_s,a_r,w_{s+1},\ldots,w_m),
		\qquad j=r+s.
		\]
		Its first $j$ vertices form exactly the set
		\[
		V_j(\pi)=A_0\cup\{w_1,\ldots,w_s\}.
		\]
		Moreover, the prescribed images of the root vertices are
		$u_i\mapsto a_i$ for every $1\le i\le r$.
		It follows that this state is good for $(T_1,e)$ if and only
		if $\{w_1,\ldots,w_s\}\in\mathcal A$. The corresponding
		equivalences hold for $(T_2,e)$ and $\mathcal B$, and for
		$(T,e)$ and $\mathcal C$. Therefore
		\[
		C_{\mathbf a}(T_1,e)=c_U(\mathcal A),\qquad
		C_{\mathbf a}(T_2,e)=c_U(\mathcal B),\qquad
		C_{\mathbf a}(T,e)=c_U(\mathcal C).
		\]
		Lemma~\ref{lem:prefix} consequently gives, for this assignment,
		\[
		C_{\mathbf a}(T_1,e)+C_{\mathbf a}(T_2,e)
		\le (m+1)!+C_{\mathbf a}(T,e).
		\]
		
		Finally, $m=n-r$ is the same for every assignment. Summing
		the preceding inequality over $\mathcal I_H$ and
		using~\eqref{eq:assignment-sum}, we obtain
		\[
		\begin{aligned}
			C(T_1,e)+C(T_2,e)
			&=
			\sum_{\mathbf a\in\mathcal I_H}
			\bigl(C_{\mathbf a}(T_1,e)+C_{\mathbf a}(T_2,e)\bigr)\\
			&\le
			\sum_{\mathbf a\in\mathcal I_H}
			\bigl((m+1)!+C_{\mathbf a}(T,e)\bigr)\\
			&=
			r!|E(H)|(m+1)!+C(T,e)\\
			&=
			M+C(T,e),
		\end{aligned}
		\]
		where the last equality follows from $m+1=n-r+1$
		and~\eqref{eq:counts}. This proves~\eqref{eq:root-gluing}.
		\iffalse
		Apply Lemma~\ref{lem:prefix} to these three families. Under
		\eqref{eq:fiber}, the three prefix counts are exactly the numbers
		of good states with the fixed assignment for $T_1$, $T_2$ and
		$T$. Summing~\eqref{eq:prefix-gluing} over all $r!|E(H)|$ possible
		assignments gives~\eqref{eq:root-gluing}, since
		$r!|E(H)|(m+1)!=M$ by~\eqref{eq:counts}.
		\fi
	\end{proof}
	
	\subsection{Completing the proof of Theorem~\ref{thm:main}}\label{subsec:induction}
	In this section we obtain an upper bound for $M$ by applying Lemmas~\ref{lem:leaf} and
	\ref{lem:root-gluing},  and use this to complete the proof of Theorem~\ref{thm:main}.
	
	\begin{theorem}\label{thm:rooted-count}
		For every tight $r$-tree $T$ with $t\ge1$ edges and every
		$e\in E(T)$,
		\begin{equation}\label{eq:rooted-count}
			M\le C(T,e)+(t-1)W.
		\end{equation}
	\end{theorem}
	
	\begin{proof}
		We proceed by induction on $t$. For $t=1$, the claim
		is~\eqref{eq:rooted-base}. Suppose $t\ge2$, and apply
		Lemma~\ref{lem:structure} at $e$. In case~\ref{case:leaf}, the smaller tight tree $S$ has $t-1$
		edges and for some edge $f \in E(S)$ the edge $e$ intersects $f$ in exactly $r - 1$ vertices. The induction hypothesis rooted at $f$, followed by
		Lemma~\ref{lem:leaf}, gives
		\[
		M\le C(S,f)+(t-2)W\le C(T,e)+(t-1)W.
		\]
		
		In case~\ref{case:gluing}, write $t_i=|E(T_i)|$. Then $t_i<t$
		for $i=1,2$ and $t_1+t_2=t+1$. Add the two induction inequalities
		and apply Lemma~\ref{lem:root-gluing} to obtain
		\[
		\begin{aligned}
			2M
			&\le C(T_1,e)+C(T_2,e)+(t_1+t_2-2)W\\
			&\le M+C(T,e)+(t-1)W.
		\end{aligned}
		\]
		Subtracting $M$ completes the induction.
	\end{proof}
	
	\begin{proof}[Proof of Theorem~\ref{thm:main}]
		If $E(H)=\varnothing$, the assertion is immediate. Otherwise
		$n\ge r$. Fix any edge $e$ of $T$. Since $H$ is $T$-free,
		$C(T,e)=0$. Theorem~\ref{thm:rooted-count} and~\eqref{eq:counts}
		give
		\[
		r!\,|E(H)|\,(n-r+1)! = M \le (t-1)W 
		= (t-1)(r-1)!\,|\sh H|\,(n-r+1)!.
		\]
		Cancelling $(r-1)!(n-r+1)!$ proves
		$r|E(H)|\le(t-1)|\sh H|$. Finally,
		$|\sh H|\le\binom{n}{r-1}$ gives~\eqref{eq:kalai-binomial}.
	\end{proof}
	
	\iffalse
	
	\section{Characterization of equality}\label{sec:equality}
	
	For an $(r-1)$-set $A\subseteq V(H)$, write
	\[
	N_H(A)=\{x\in V(H)\setminus A:A\cup\{x\}\in E(H)\},
	\qquad d_H(A)=|N_H(A)|.
	\]
	Thus $A\in\sh H$ if and only if $d_H(A)>0$. We call a tight
	$r$-tree with $t$ edges a \emph{sunflower} if its edges are
	\[
	F\cup\{z_1\},\ldots,F\cup\{z_t\},
	\qquad |F|=r-1,
	\]
	where the $z_i$ are distinct vertices outside $F$.
	
	\begin{theorem}[Equality in the shadow bound]\label{thm:equality-shadow}
		Let $r\ge2$ and $t\ge2$, let $T$ be a tight $r$-tree with $t$
		edges, and put $b=t+r-2$. If $H$ is a $T$-free $r$-graph, then
		\begin{equation}\label{eq:equality-shadow}
			r|E(H)|=(t-1)|\sh H|
		\end{equation}
		if and only if the applicable one of the following conditions holds.
		\begin{enumerate}[label=\textup{(\roman*)},leftmargin=*]
			\item If $T$ is a sunflower, then
			\[
			d_H(A)=t-1\qquad\text{for every }A\in\sh H.
			\]
			\item If $T$ is not a sunflower, then there is a family
			$\mathcal B\subseteq\binom{V(H)}b$ such that
			\begin{equation}\label{eq:equality-blocks}
				E(H)=\bigcup_{B\in\mathcal B}\binom Br,
				\qquad |B\cap B'|\le r-2
				\quad\text{for distinct }B,B'\in\mathcal B.
			\end{equation}
		\end{enumerate}
		Vertices outside the union of the blocks are allowed and are isolated.
	\end{theorem}
	
	\begin{proof}
		The assertion is immediate when $E(H)=\varnothing$, with
		$\mathcal B=\varnothing$ in the second case. Assume henceforth
		that $E(H)\ne\varnothing$ and that~\eqref{eq:equality-shadow}
		holds. With the notation of~\eqref{eq:counts}, equality says
		\begin{equation}\label{eq:equality-MW}
			M=(t-1)W.
		\end{equation}
		
		\medskip\noindent
		\emph{Equality forces constant positive codegree.}
		Fix a tight-tree construction of $T$, and write its last edge as
		$e=F\cup\{\ell\}$, where $\ell$ is the newly introduced vertex.
		Delete $e$ and $\ell$ to obtain a tight tree $S$ with $t-1$ edges,
		and choose an edge $f=F\cup\{u\}$ of $S$. Order $f$ with the
		vertices of $F$ first and $u$ last. Theorem~\ref{thm:rooted-count}
		and~\eqref{eq:equality-MW} give
		\begin{equation}\label{eq:equality-rooted-lower}
			C(S,f)\ge M-(t-2)W=W.
		\end{equation}
		
		For each $\pi\in\Wcal$, there is at most one good state for
		$(S,f)$. Indeed, if $i<j$ are two good positions, choose an
		embedding witnessing the state at $i$. Its image lies in the
		first $i$ entries, and it maps $F$ to the first $r-1$ entries.
		The vertex $v_j$ is unused and completes these entries to an
		edge of $H$. Sending $\ell$ to $v_j$ therefore extends the
		embedding to $T$, a contradiction. Together with
		\eqref{eq:equality-rooted-lower}, this proves that every
		$\pi\in\Wcal$ has exactly one good state for $(S,f)$.
		
		\bigskip
		
		Suppose that $A\in\sh H$ satisfies $d_H(A)\le t-2$. Choose
		a permutation beginning with the vertices of $A$, followed by
		all vertices of $N_H(A)$, followed by the remaining vertices.
		Every marked position $j$ in this permutation satisfies
		\[
		j\le r-1+d_H(A)\le t+r-3.
		\]
		Since $S$ has $t+r-2$ vertices, none of these positions can be
		good for $(S,f)$. This is a contradiction. Hence every member
		of $\sh H$ has codegree at least $t-1$. Double counting gives
		\[
		\sum_{A\in\sh H}d_H(A)=r|E(H)|=(t-1)|\sh H|,
		\]
		so in fact
		\begin{equation}\label{eq:equality-codegree}
			d_H(A)=t-1\qquad\text{for every }A\in\sh H.
		\end{equation}
		
		If $T$ is a sunflower, this proves necessity. Conversely,
		\eqref{eq:equality-codegree} forbids that sunflower and gives
		\eqref{eq:equality-shadow} by double counting. This proves
		the first case of the theorem.
		
		\medskip\noindent
		\emph{For a nonsunflower, every codimension-two link is a union
			of cliques.}
		Assume now that $T$ is not a sunflower. Use the construction
		already chosen to form the parent tree $D$ on $E(T)$: join each
		edge other than the first to its chosen parent. Adjacent nodes
		meet in exactly $r-1$ vertices, and for each vertex $v$ of $T$
		the nodes containing $v$ form a connected subtree. These are
		the properties established in the proof of
		Lemma~\ref{lem:structure}. The node $e$ is a leaf, and $D-e$
		represents $S$.
		
		\bigskip
		
		The nodes of $D-e$ containing all of $F$ form a nonempty
		connected subtree: they are the intersection of the subtrees
		corresponding to the vertices of $F$. This subtree is proper,
		since otherwise every edge of $T$ would contain $F$ and $T$
		would be a sunflower. There are therefore adjacent nodes
		$f_0,g_0$ in $D-e$ with $F\subseteq f_0$ and
		$F\nsubseteq g_0$. They have the form
		\begin{equation}\label{eq:equality-two-edges}
			f_0=F\cup\{u_0\},
			\qquad
			g_0=(F\setminus\{a\})\cup\{u_0,w\},
		\end{equation}
		where $a\in F$ and $w\notin f_0$.
		Rooting $D-e$ at $f_0$ and placing $g_0$ second gives a
		tight-tree construction of $S$ beginning with $f_0,g_0$.
		Indeed, in any order with parents preceding children, the
		vertex of a child edge outside its parent cannot have occurred
		earlier: otherwise the connectedness of its occurrence nodes
		would force it to belong to the parent.
		
		\bigskip
		
		For an $(r-2)$-set $Q\subseteq V(H)$, let $L_H(Q)$ be the
		graph on $V(H)\setminus Q$ with
		\[
		xy\in E(L_H(Q))
		\quad\Longleftrightarrow\quad
		Q\cup\{x,y\}\in E(H).
		\]
		Suppose that $x,y,z$ induce a three-vertex path in $L_H(Q)$.
		Thus
		\begin{equation}\label{eq:equality-link-path}
			Q\cup\{x,y\},\ Q\cup\{y,z\}\in E(H),
			\qquad Q\cup\{x,z\}\notin E(H).
		\end{equation}
		Embed the two edges in~\eqref{eq:equality-two-edges} by an
		arbitrary bijection from $F\setminus\{a\}$ to $Q$ and by
		\[
		a\mapsto x,\qquad u_0\mapsto y,\qquad w\mapsto z.
		\]
		Extend this embedding greedily along the construction of $S$
		just described. After $q$ edges have been embedded, there are
		$q+r-1$ used vertices. The attaching $(r-1)$-face of the next
		edge belongs to an already embedded edge, so its image lies in
		$\sh H$ and has codegree $t-1$ by~\eqref{eq:equality-codegree}.
		Exactly $q\le t-2$ used vertices lie outside that face. An
		unused extension therefore exists at every step.
		
		\bigskip
		
		Let $\varphi$ be the resulting embedding of $S$, and set
		$A=\varphi(F)=Q\cup\{x\}$. The image of $S$ contains exactly
		$t-1$ vertices outside $A$. One of them is $z=\varphi(w)$,
		which is not in $N_H(A)$ by~\eqref{eq:equality-link-path}.
		Thus at most $t-2$ vertices of $N_H(A)$ are used. Since
		$d_H(A)=t-1$, an unused neighbour of $A$ remains; sending
		$\ell$ to that vertex produces a copy of $T$. This is a
		contradiction.
		
		\bigskip
		
		It follows that every $L_H(Q)$ contains no induced
		three-vertex path. Each of its connected components is
		therefore complete: a shortest path between two nonadjacent
		vertices would otherwise begin with an induced three-vertex
		path. Moreover, every vertex of positive degree in $L_H(Q)$
		has degree $t-1$, since that degree is
		$d_H(Q\cup\{x\})$ for the corresponding vertex $x$.
		Consequently every nontrivial component of every $L_H(Q)$
		is a copy of $K_t$.
		
		\medskip\noindent
		\emph{The local cliques determine complete blocks.}
		For $A\in\sh H$, put
		\[
		B_A=A\cup N_H(A).
		\]
		By~\eqref{eq:equality-codegree}, $|B_A|=t+r-2=b$.
		Choose $a\in A$ and $x\in N_H(A)$, and let
		\[
		Q=A\setminus\{a\},\qquad A'=Q\cup\{x\}.
		\]
		The component of $L_H(Q)$ containing $a$ is a clique with
		vertex set $\{a\}\cup N_H(A)$. The vertex $x$ belongs to
		this clique, so
		\[
		N_H(A')=
		\bigl(\{a\}\cup N_H(A)\bigr)\setminus\{x\}.
		\]
		In particular, $A'\in\sh H$ and
		\begin{equation}\label{eq:equality-block-exchange}
			B_{A'}=B_A.
		\end{equation}
		
		Every $(r-1)$-subset of $B_A$ can be reached from $A$ by
		successive single-vertex exchanges. At each step, the entering
		vertex belongs to the current block but not to the current
		face, and hence belongs to the neighbourhood of that face.
		Repeated application of~\eqref{eq:equality-block-exchange}
		therefore shows that
		\begin{equation}\label{eq:equality-all-faces}
			A'\in\sh H,\qquad B_{A'}=B_A
			\quad\text{for every }A'\in\binom{B_A}{r-1}.
		\end{equation}
		It follows that every $r$-subset of $B_A$ is an edge of $H$.
		Conversely, every edge of $H$ is contained in $B_A$ for each
		of its $(r-1)$-faces $A$.
		
		Let $\mathcal B$ be the family of distinct sets $B_A$ as
		$A$ ranges over $\sh H$. The preceding paragraph proves
		\[
		E(H)=\bigcup_{B\in\mathcal B}\binom Br.
		\]
		If two members $B,B'$ shared an $(r-1)$-set $A$, then
		\eqref{eq:equality-all-faces} would give $B=B_A=B'$.
		Thus distinct blocks intersect in at most $r-2$ vertices,
		as required in~\eqref{eq:equality-blocks}.
		
		\medskip\noindent
		\emph{Conversely, the block construction attains equality.}
		Suppose $H$ has the form~\eqref{eq:equality-blocks}. Each
		edge belongs to a unique block, and edges belonging to
		different blocks cannot meet in $r-1$ vertices. Following
		the parent edges in a tight-tree construction, any copy of
		$T$ in $H$ would therefore lie in a single block. This is
		impossible, since each block has $t+r-2$ vertices and $T$
		has $t+r-1$ vertices.
		
		Each member of $\sh H$ belongs to exactly one block, and
		its neighbourhood consists of the $b-(r-1)=t-1$ remaining
		vertices of that block. Thus~\eqref{eq:equality-codegree}
		holds, and double counting proves~\eqref{eq:equality-shadow}.
	\end{proof}
	
	\begin{corollary}[Equality in the binomial bound]
		\label{cor:equality-binomial}
		Let $r\ge2$, $t\ge2$, $n\ge r$, and let $T$ be a tight
		$r$-tree with $t$ edges. An $n$-vertex $T$-free $r$-graph $H$
		satisfies
		\[
		|E(H)|=\frac{t-1}{r}\binom n{r-1}
		\]
		if and only if the applicable one of the following conditions holds.
		\begin{enumerate}[label=\textup{(\roman*)},leftmargin=*]
			\item If $T$ is a sunflower, every $(r-1)$-subset of $V(H)$
			has codegree $t-1$; equivalently, $H$ is a simple
			$(r-1)$-$(n,r,t-1)$ design.
			\item If $T$ is not a sunflower, there is a Steiner system
			$S(r-1,t+r-2,n)$ with block family $\mathcal B$ such that
			\[
			E(H)=\bigcup_{B\in\mathcal B}\binom Br.
			\]
		\end{enumerate}
	\end{corollary}
	
	\begin{proof}
		Since $t-1>0$, equality in
		\[
		|E(H)|\le\frac{t-1}{r}|\sh H|
		\le\frac{t-1}{r}\binom n{r-1}
		\]
		holds if and only if equality holds in the shadow bound and
		$\sh H=\binom{V(H)}{r-1}$. Apply
		Theorem~\ref{thm:equality-shadow}. In the second case, the
		blocks must cover every $(r-1)$-set exactly once, which is
		precisely the defining property of the stated Steiner system.
	\end{proof}
	
	\fi
	\begin{thebibliography}{100}
		
		\bibitem{ASS} M.~Ajtai, M.~Simonovits, E.~Szemerédi, \emph{Solution of the Erd\H{o}s-S\'{o}s Conjecture}, Presentation at Bondy Conference, Montreal, 2003.
		
		\bibitem{BPSS2019} G.~Besomi, M.~Pavez-Sign{\'e}, and M.~Stein, \emph{Degree conditions for embedding trees}, SIAM J. Discrete Math. \textbf{33} (2019), 1521--1555. 
		
		\bibitem{BPSS2021} G.~Besomi, M.~Pavez-Sign{\'e}, and M.~Stein, \emph{On the Erd{\H{o}}s--S{\'o}s conjecture for trees with bounded degree}, Combin. Probab. Comput. \textbf{30} (2021), 741--761. 
		
		\bibitem{BD} S.~Brandt and E.~Dobson, \emph{The Erd{\H{o}}s--S{\'o}s conjecture for graphs of girth 5}, Discrete Math. \textbf{150} (1996), 411--414. 
		
		\bibitem{DPRS} A.~Davoodi, D.~Piguet, H.~\v{R}ada, and N.~Sanhueza-Matamala, \emph{The asymptotic version of the Erd{\H{o}}s--S{\'o}s conjecture and beyond}, arXiv:2603.17755, 2026. 
		
		
		\bibitem{Epoch}
		T.~Adamczewski and T.~F.~Bloom,
		\emph{FrontierMath Erd\H{o}s}, Epoch AI report, September 2026.
		\url{https://epoch.ai/files/frontiermath-erdos.pdf}.
		
		\bibitem{Erdos}
		P.~Erd\H{o}s,
		\emph{Extremal problems in graph theory},
		in: Theory of Graphs and its Applications (M.~Fiedler, ed.),
		Academic Press, New York, 1965, 29--36.
		
		\bibitem{Erdos1964} P.~Erd{\H{o}}s, \emph{Extremal problems in graph theory}, in: Theory of Graphs and its Applications (Proc. Sympos. Smolenice, 1963), Publ. House Czechoslovak Acad. Sci., Prague, 1964, pp.~29--36. 
		
		\bibitem{EG}
		P.~Erd\H{o}s and T.~Gallai,
		\emph{On maximal paths and circuits of graphs},
		Acta Math. Acad. Sci. Hungar. \textbf{10} (1959), 337--356.
		
		\bibitem{ErdosSos1970} P.~Erd{\H{o}}s and V.~T.~S{\'o}s, \emph{Some remarks on Ramsey's and Tur{\'a}n's theorem}, in: Combinatorial Theory and its Applications, II (Proc. Colloq., Balatonf{\"u}red, 1969), Colloq. Math. Soc. J{\'a}nos Bolyai, Vol.~4, North-Holland, Amsterdam, 1970, pp.~395--404.
		
		\bibitem{FanSun2007} G.~Fan and L.~Sun, \emph{The Erd{\H{o}}s--S{\'o}s conjecture for spiders}, Discrete Math. \textbf{307} (2007), 3055--3062. 
		
		
		
		\bibitem{FF}
		P.~Frankl and Z.~F\"uredi,
		\emph{Exact solution of some Tur\'an-type problems}, J. Combin. Theory Ser. A \textbf{45} (1987), no.~2, 226--262.
		
		\bibitem{FJKMVtrunk}
		Z.~F\"uredi, T.~Jiang, A.~Kostochka, D.~Mubayi and J.~Verstra\"ete,
		\emph{Hypergraphs not containing a tight tree with a bounded trunk},
		SIAM J. Discrete Math. \textbf{33} (2019), no.~2, 862--873.
		\href{https://arxiv.org/abs/1712.04081}{arXiv:1712.04081}.
		
		\bibitem{FJKMVpaths}
		Z.~F\"uredi, T.~Jiang, A.~Kostochka, D.~Mubayi and J.~Verstra\"ete,
		\emph{Tight paths in convex geometric hypergraphs},
		Adv. Combin. \textbf{2020} (2020), Paper No.~1, 14~pp.
		\href{https://doi.org/10.19086/aic.12044}{doi:10.19086/aic.12044}.
		
		\bibitem{GKLO} S.~Glock, D.~K\"{u}hn, A.~Lo, D.~Osthus, \emph{The existence of designs via iterative absorption: hypergraph $F$-designs for arbitrary $F$}, Memoirs of the American Mathematical Society \textbf{284} (2023), monograph 1406).
		
		
		\bibitem{GoerlichZak2016} A.~Goerlich and A.~\.{Z}ak, \emph{On Erd{\H{o}}s--S{\'o}s conjecture for trees of large size}, Electron. J. Combin. \textbf{23} (2016), Paper~P1.52. 
		
		\bibitem{HRSW2020} F.~Havet, B.~Reed, M.~Stein, and D.~R.~Wood, \emph{A variant of the Erd{\H{o}}s--S{\'o}s conjecture}, J. Graph Theory \textbf{94} (2020), 131--158.
		
		
		\bibitem{Keevash}
		P.~Keevash,
		\emph{The existence of designs},
		\href{https://arxiv.org/abs/1401.3665}{arXiv:1401.3665}, 2014;
		revised 2024.
		
		\bibitem{KupitzPerles} Y. S. Kupitz, M. Perles, \emph{Extremal theory for convex matchings in convex geometric graphs}, Discrete Comput. Geom. \textbf{15}, (1996), 195--220. 
		
		
		\bibitem{McLennan2005} A.~McLennan, \emph{The Erd{\H{o}}s--S{\'o}s conjecture for trees of diameter four}, J. Graph Theory \textbf{49} (2005), 291--301.
		
		\bibitem{MubayiV2016} D.~Mubayi and J.~Verstra\"ete, \emph{Counting trees in graphs}, Electron. J. Combin. \textbf{23} (2016), no.~3, Paper P3.39.
		
		\bibitem{Pokrovskiy2024} A.~Pokrovskiy, \emph{Hyperstability in the Erd{\H{o}}s--S{\'o}s conjecture}, arXiv:2409.15191, 2024. 
		
		\bibitem{ReedStein2026} B.~Reed and M.~Stein, \emph{The Erd{\H{o}}s--S{\'o}s conjecture in dense graphs}, arXiv:2609.05417, 2026. 
		
		
		\bibitem{Rodl}
		V.~R\"odl,
		\emph{On a packing and covering problem},
		European J. Combin. \textbf{6} (1985), no.~1, 69--78.
		
		\bibitem{Rozhon2019} V.~Rozho\v{n}, \emph{A local approach to the Erd{\H{o}}s--S{\'o}s conjecture}, SIAM J. Discrete Math. \textbf{33} (2019), 643--664.
		
		\bibitem{SacleWozniak1997} J.-F.~Sacl{\'e} and M.~Wo{\'z}niak, \emph{The Erd{\H{o}}s--S{\'o}s conjecture for graphs without $C_4$}, J. Combin. Theory Ser. B \textbf{70} (1997), 367--372. 
		
		\bibitem{Stein} M.~Stein, \emph{Tree containment and degree conditions}, in: A.~M.~Raigorodskii and M.~T.~Rassias (eds.), Discrete Mathematics and Applications, Springer Optimization and Its Applications, Springer, 2020, pp.~459--486. 
		
		\bibitem{TinerTomlin2022} G.~Tiner and Z.~Tomlin, \emph{On the Erd{\H{o}}s--S{\'o}s conjecture for $k=9$}, Alabama J. Math. \textbf{45} (2022), 37--45.
		
		\bibitem{Wilson2025} C.~Wilson, \emph{A Tight Lower bound on Trees in Graphs}, https://arxiv.org/abs/2512.14890. 
		
		\bibitem{Wozniak1996} M.~Wo{\'z}niak, \emph{On the Erd{\H{o}}s--S{\'o}s conjecture}, J. Graph Theory \textbf{21} (1996), 229--234.
		
	\end{thebibliography}
	
\end{document}
